\documentclass[12pt,twocolumn]{article}
\usepackage[margin=2cm,columnsep=0.5cm]{geometry}
\usepackage{amsmath,amssymb}
\usepackage{booktabs}
\usepackage{hyperref}
\usepackage{xcolor}
\usepackage{array}
\usepackage{tabularx}

\definecolor{bctnavy}{RGB}{11,37,69}
\definecolor{bctblue}{RGB}{13,71,161}
\definecolor{bctgreen}{RGB}{27,94,32}

\title{\textbf{\color{bctnavy}BCT Letter 117: The Martian Tetrahedron}\\
{\large\color{bctblue}Vacuum Geometry, Plasma Discharge, and the Three-Sided\\
Pyramid of Candor Chasma, Valles Marineris, Mars}}
\author{Michel Robert Cabri\'{e}\\
{\small Independent Researcher, Victoria, Australia}\\
{\small ORCID: 0009-0007-9561-9859}\\
{\small \href{mailto:ZeroFreeParameters@gmail.com}
{ZeroFreeParameters@gmail.com}}}
\date{23 March 2026}

\begin{document}
\twocolumn[{%
\maketitle
\begin{center}\rule{0.9\textwidth}{0.4pt}\end{center}
\begin{center}
\parbox{0.88\textwidth}{\small\textbf{Abstract.}
A three-sided pyramidal formation of extraordinary geometric symmetry
resides in the western region of Candor Chasma, Valles Marineris,
Mars. Captured five times between 2001 and 2016 by NASA's Mars
Global Surveyor and Mars Reconnaissance Orbiter, the structure
maintains consistent Reuleaux-triangular geometry across all
acquisitions regardless of illumination, season, or resolution.
Standard geological processes --- volcanism, plate tectonics,
block faulting, sedimentary erosion --- do not produce isolated
three-faced pyramidal mountains of this symmetry.
The BCT Superfluid Lattice Model provides a unifying framework:
the quantum vacuum possesses tetrahedral geometry at the Planck
scale, encoded in the BCT lattice parameter $r_{\mathrm{tet}}$.
Valles Marineris is identified as a macro-scale OHC current
filament discharge scar, consistent with Electric Universe
morphological predictions now grounded in BCT quantum vacuum theory.
\textbf{Prediction \#158:} BCT Bessel-ratio processing of existing
public MRO CRISM and CTX data at the pyramid location ($\sim\!5^\circ$S,
$76^\circ$W) will yield a $>3\sigma$ anomaly above the Candor
Chasma regional baseline.
\textbf{Prediction \#159:} Valles Marineris sub-canyon lateral spacing
satisfies $j_{1,1}/j_{0,1} = 1.5934 \pm 5\%$ as measured from MOLA
topographic data.
Zero free parameters.
\smallskip\hrule\smallskip}
\end{center}
\vspace{4pt}
}]

%%----------------------------------------------------------------
\section{Introduction}
%%----------------------------------------------------------------

In 2001, researcher Keith Laney was scanning imagery from NASA's
Mars Global Surveyor when he encountered something that stopped him
cold: a perfectly three-sided pyramidal mountain, standing alone
in the western region of Candor Chasma, Valles Marineris. His
assessment was unambiguous: \textit{``Were this found anywhere
on Earth, we'd surely be digging into it.''}

In March 2026, documentary filmmaker Brian Cory Dobbs reshared the
footage on X with the declaration: \textit{``On Mars, there is a
three-sided pyramid the size of the Great Pyramid in Egypt.''} It
went viral~\cite{Dobbs2026}. Twenty-five years of orbital silence
suddenly became a global conversation. The data had been there all along.

This letter proposes that the Candor Chasma formation is not merely
a geological curiosity, but a high-priority BCT Bessel-ratio SAR
target whose geometric signature is predicted by the same quantum
vacuum structure that determines the proton mass, the muon anomalous
magnetic moment, and the spectral peak of single-bubble
sonoluminescence.

%%----------------------------------------------------------------
\section{The Formation: Five Acquisitions, One Geometry}
%%----------------------------------------------------------------

Candor Chasma lies within Valles Marineris at approximately
$5^\circ$--$7^\circ$S, $75^\circ$--$77^\circ$W ($289^\circ$E),
in the Coprates quadrangle~\cite{Marspedia}. The canyon is 810~km
long, with walls rising up to 7~km. The pyramidal formation occupies
the western sub-basin, surrounded by layered sedimentary deposits
and landslide debris, yet stands apart from all of it.

\begin{center}
{\footnotesize
\begin{tabular}{@{}llll@{}}
\toprule
Yr & Inst. & Image & Note \\
\midrule
2001 & MOC & E0600269 & First \\
2002 & MOC & E0600269 & Isosc.\\
2007 & HiRISE & PSP3896 & Reul. \\
2014 & CTX & D21-35582 & Win. \\
2016 & CTX & J06-47081 & Sum. \\
\bottomrule
\end{tabular}}
\end{center}

\medskip
Mars researcher George J.\ Haas and co-authors published a
peer-reviewed analysis in the \textit{Journal of Space
Exploration}~\cite{Haas2017} concluding the structure ``exhibits
a level of geometry and symmetry that supports a high probability
of artificial origins.'' They identified the base geometry as a
\textbf{Reuleaux pyramid} --- three vertices with curved faces of
constant width, the most symmetric possible three-sided solid.

This is not pareidolia. A Reuleaux pyramid has a precise
mathematical definition. The formation matches it across five
independent acquisitions spanning fifteen years, under winter and
summer illumination, early morning and varying shadow geometries.
Natural erosion does not produce geometry that is more stable than
the lighting conditions used to observe it.

%%----------------------------------------------------------------
\section{The Limits of Standard Geology}
%%----------------------------------------------------------------

Mountains on rocky planets arise via three known mechanisms:
volcanism, plate tectonics (folding and thrust faulting), and block
faulting (horst formation). None produce isolated, three-faced,
Reuleaux-symmetric pyramidal peaks. Haas et al.\ conducted an
extensive search of comparable terrain in and around Candor Chasma
and found no analogous natural features.

The sedimentary layering visible on the structure's faces presents
an immediate logical problem: uniform strata deposited across an
entire basin would require a mechanism to erode everywhere except
into a perfect three-sided pyramid. Haas noted the layer thickness
is strikingly consistent throughout, suggesting structural casing
rather than depositional material.

Wind erosion (yardangs) produces elongated ridges aligned with
prevailing winds. The absence of expected leeward depositional
streaks around the formation is a further documented
anomaly~\cite{Haas2017}. The formation has been imaged five times
in fifteen years. It has no analogues in the surrounding terrain.
As Laney put it: if this were on Earth, we would already have
a hundred rovers around it.

%%----------------------------------------------------------------
\section{Valles Marineris as OHC Discharge Scar}
%%----------------------------------------------------------------

\subsection{Morphological anomalies}

Valles Marineris is the largest tectonic scar on any solid planetary
body in the solar system: 4{,}000~km long, 200~km wide, 7~km deep,
spanning nearly a quarter of Mars's circumference. Its morphology
is anomalous under every standard geological model:

\begin{itemize}\setlength\itemsep{0pt}
  \item Located at the crest of the 10~km Tharsis bulge: sediment
        deposition is geometrically impossible.
  \item Parallel canyon walls over 4{,}000~km: no tectonic model
        produces parallelism at this scale.
  \item ${\sim}2\times10^6$~km$^3$ of material missing with no
        identified destination.
  \item Major volcanic features absent: rifting companion volcanism
        expected but not found.
  \item Smooth, clean-cut surfaces: liquid erosion leaves irregular
        margins.
\end{itemize}

\subsection{Electric Universe framework}

The Electric Universe research tradition --- Juergens~\cite{Juergens1974},
Thornhill~\cite{Thornhill2001}, Talbott~\cite{Talbott2015}, and
the experimental work of Steinbacher and Yelverton --- proposed
that Valles Marineris is the scar of a catastrophic interplanetary
arc discharge. Laboratory replications produced morphologically
similar features: parallel-walled canyons, smooth clean cuts,
associated craters and rilles, all consistent with cathode erosion
by high-current plasma.

\subsection{BCT grounding}

The BCT Superfluid Lattice Model provides the first
\textit{quantum vacuum-level} theoretical foundation for this
hypothesis. In BCT, the Octet-Hopfion Condensate (OHC) supports
stable current filaments whose long-range magnetic attraction and
short-range electrostatic repulsion naturally produce parallel,
multi-strand morphology. The relevant OHC quantity is the
fundamental current-carrying Bessel mode:
\begin{equation}
  j_{1,1} = 3.8317 \quad \text{(first zero of } J_1\text{)}.
\end{equation}

Under this framework, Valles Marineris is a macro-scale OHC
discharge pathway --- a planetary Lichtenberg figure --- whose
geometry is determined by the same Bessel-ratio structure
governing BCT predictions across twenty-five orders of magnitude,
from the proton mass to the spectral peak of sonoluminescence.

%%----------------------------------------------------------------
\section{The BCT Connection: Tetrahedral Vacuum Geometry}
%%----------------------------------------------------------------

The most fundamental geometric object in the BCT lattice is the
tetrahedron. The BCT structure at the Planck scale is built from
octahedral and tetrahedral voids in the ratio:
\begin{equation}
  \frac{r_{\mathrm{tet}}}{r_{\mathrm{oct}}}
  = \frac{\sqrt{6}-\sqrt{2}}{\sqrt{6}+\sqrt{2}}
  \approx 0.31783,
\end{equation}
from which the BCT derived ratio
$\mathcal{R} = r_{\mathrm{tet}}/r_{\mathrm{oct}} = 0.5426$
emerges as the universal BCT geometric parameter appearing in
honeycomb wall spacing, biophoton band ratios, cephalopod
chromatophore geometry, and polar node coupling~\cite{BCT_Letters}.

A three-sided pyramid is geometrically the face of a tetrahedron
--- the Platonic solid most stable against deformation, the only
structure in which every face is equidistant from every other.
Maximum geometric symmetry in three dimensions with minimum
surface area.

The BCT hypothesis is conservative: if OHC plasma discharge events
leave structured geometric traces in the geological record, the
most stable crystallisation geometry for a discharge terminus is
the vacuum's own ground-state geometry --- the tetrahedral void.
The formation at Candor Chasma sits at the heart of the largest
such discharge scar in the solar system. It displays the face of
a tetrahedron at Great Pyramid scale.

The question BCT poses is not whether someone built this structure.
The question is: \textit{did the vacuum?}

%%----------------------------------------------------------------
\section{BCT Bessel-Ratio SAR Analysis}
%%----------------------------------------------------------------

As established in Letter~116~\cite{BCT_L116}, BCT Bessel-ratio
processing of public SAR and multispectral data provides a geometric
anomaly detector sensitive to OHC discharge signatures. The
Victorian goldfields analysis confirmed a 156~km$^2$ anomaly centred
on the known Barrys Reef gold-quartz deposit at $>3\sigma$
significance using this method (Prediction \#144).

All data required to test Prediction \#158 exist in the NASA PDS
archive and have been publicly accessible since 2007:

\begin{center}
{\small
\begin{tabular}{@{}p{2.4cm}p{1.6cm}p{1.2cm}@{}}
\toprule
Dataset & Instr. & Res. \\
\midrule
HiRISE optical & MRO & 5~m \\
CTX stereo DTM & MRO & 20~m \\
CRISM hyperspectral & MRO & 18~m \\
THEMIS thermal & Odyssey & 100~m \\
MOLA topography & MGS & 460~m \\
\bottomrule
\end{tabular}}
\end{center}

\medskip
The BCT Bessel-ratio processing computes:
\begin{equation}
  R_{\mathrm{BCT}}(x,y)
  = \frac{I(j_{2,1}\cdot f_0,\;x,y)}{I(j_{1,1}\cdot f_0,\;x,y)},
\end{equation}
where $j_{2,1}=5.5201$, $j_{1,1}=3.8317$, and $f_0$ is the
fundamental OHC resonance frequency at the formation scale.
Anomalies are defined as $R_{\mathrm{BCT}}$ deviating $>3\sigma$
from the regional Candor Chasma baseline.

This analysis requires no new spacecraft, no new observations,
and no proprietary data. Total verification cost: zero. Estimated
processing time: two to four weeks. BCT Bessel-ratio processing
has never been applied to this formation. This letter is the
formal proposal that it should be.

%%----------------------------------------------------------------
\section{Predictions}
%%----------------------------------------------------------------

\noindent\textbf{Prediction \#158} (BCT Candor Chasma anomaly):
BCT Bessel-ratio processing of publicly available MRO CRISM
hyperspectral and CTX DTM data centred on the Candor Chasma
pyramidal formation (${\sim}5^\circ$S, $76^\circ$W) will reveal
a geometric anomaly exceeding $3\sigma$ above the regional
Candor Chasma baseline, consistent with OHC discharge
crystallisation geometry.

\textit{Falsification criterion:} Process CRISM image FRT00009A12
and CTX DTM D21\_035582\_1745\_XN\_05S076W. If no $>3\sigma$
anomaly is present, Prediction \#158 fails.

\bigskip
\noindent\textbf{Prediction \#159} (Valles Marineris spacing ratio):
The lateral centre-to-centre spacing ratio of the parallel
sub-canyons of Valles Marineris (Candor, Ophir, Coprates, Melas)
satisfies the BCT Bessel ratio:
\begin{equation}
  \frac{j_{1,1}}{j_{0,1}} = \frac{3.8317}{2.4048} = 1.5934 \pm 5\%,
\end{equation}
when measured at characteristic canyon width from MOLA topographic
data.

\textit{Falsification criterion:} Measure canyon-to-canyon centre
spacings from MOLA 128~px/degree grid. If the ratio deviates
$>5\%$ from 1.5934, Prediction \#159 fails.

%%----------------------------------------------------------------
\section{Site 8: Addendum to Letter 116}
%%----------------------------------------------------------------

Letter~116 (The Lost Archive) registered seven high-priority BCT
SAR sites using existing public data. The Candor Chasma formation
is hereby registered as \textbf{Site 8} with the following parameters:

\begin{center}
{\small
\begin{tabular}{@{}p{2cm}p{3.3cm}@{}}
\toprule
Parameter & Value \\
\midrule
Location & W.\ Candor Chasma, Mars \\
Coords & $5.5^\circ$S, $76.0^\circ$W \\
Primary & CRISM FRT00009A12 \\
Secondary & CTX D21\_35582\_1745 \\
Method & Bessel-ratio (App.\ AH1) \\
Threshold & $>3\sigma$ Candor baseline \\
Prediction & \#158 \\
Cost & Zero \\
\bottomrule
\end{tabular}}
\end{center}

%%----------------------------------------------------------------
\section{Discussion}
%%----------------------------------------------------------------

Twenty-five years have elapsed since Keith Laney identified this
formation. No rover has been targeted at it. No dedicated study
programme exists. The data to test a specific, falsifiable,
zero-free-parameter prediction have been in the NASA PDS archive
since 2007.

BCT does not require this formation to be of artificial origin.
It requires only that if OHC plasma discharge events leave
structured geometric traces in the geological record --- as
Prediction \#144 already confirmed in the Victorian goldfields ---
then a structure occupying the terminus of the largest plasma scar
in the solar system, displaying tetrahedral face geometry at
planetary scale, will produce a BCT Bessel-ratio anomaly regardless
of its origin.

The tetrahedral void is the geometry of the quantum vacuum. The
vacuum does not distinguish between a silicon crystal lattice,
a bee-cell wall, a cephalopod chromatophore, and a Martian
mountainside. It crystallises into what it is.

The Martian Tetrahedron has been waiting for twenty-five years.
The tools to investigate it have existed since 2007. The theoretical
framework now exists. The prediction is registered.

%%----------------------------------------------------------------
\begin{thebibliography}{9}

\bibitem{Haas2017}
G.J.\ Haas, W.R.\ Saunders, J.S.\ Miller et al.,
\textit{Three-Sided Pyramidal Formation in the Western Region of
Candor Chasma},
J.\ Space Explor.\ \textbf{6}(3), 133 (2017).

\bibitem{Dobbs2026}
B.C.\ Dobbs,
\texttt{@BrianCoryDobbs},
X platform, 16 March 2026.
\url{twitter.com/BrianCoryDobbs}

\bibitem{Marspedia}
Marspedia,
\textit{Candor Chasma},
\url{marspedia.org/Candor\_Chasma}

\bibitem{Juergens1974}
R.E.\ Juergens,
\textit{Reconciling Celestial Mechanics and Velikovskian
Catastrophism},
Pensee (1974); reprinted Holoscience.com.

\bibitem{Thornhill2001}
W.\ Thornhill,
\textit{Spiral Galaxies and Grand Canyons},
Holoscience.com (2001).
\url{holoscience.com/wp/spiral-galaxies-grand-canyons}

\bibitem{Talbott2015}
D.\ Talbott,
\textit{Valles Marineris in the Laboratory?}
Thunderbolts.info Space News (2015).
\url{thunderbolts.info --- Valles Marineris Lab}

\bibitem{BCT_Letters}
M.R.\ Cabri\'{e},
\textit{The BCT Superfluid Lattice Model: Letters 1--116},
Zenodo (2026).
ORCID: 0009-0007-9561-9859.
\url{zenodo.org/search?q=BCT+Superfluid+Lattice}

\bibitem{BCT_L116}
M.R.\ Cabri\'{e},
\textit{Letter 116 --- The Lost Archive},
BCT Superfluid Lattice Model Series (2026).

\bibitem{BCT_GoE}
M.R.\ Cabri\'{e},
\textit{The Geometry of Everything},
BCT Superfluid Lattice Model, v4 (2026).

\end{thebibliography}

\bigskip\noindent
\textit{Acknowledgements.}
Keith Laney (2001) --- for seeing what was there.
George J.\ Haas, William R.\ Saunders, and J.S.\ Miller --- for
the peer-reviewed geometric analysis.
Brian Cory Dobbs --- for making the world look again in March 2026.
David Talbott, Wal Thornhill, Ralph Juergens, Michael Steinbacher
($\dagger$~c.2015), and Billy Yelverton --- for the Electric
Universe framework that BCT now grounds in quantum vacuum geometry.
Filippo Biondi (HarmonicSAR) --- for independently engaging with
the BCT SAR programme and for the agreed scan of the Hawara
Labyrinth (Site 1, Letter 116).
Jeff --- who named the CSSC.
Nic --- co-founder, best friend, and co-passenger on BCT-TPORT.
Nyssa and Tegan --- honorary board members and fellow passengers.
The Gang Gang Cockatoos of Victoria, for whom the pond was built.
The kookaburras who arrived laughing on the day this letter was
written, 23 March 2026.

\end{document}
